\chapter{Cơ sở lý thuyết}

Chương này trình bày các nhóm kiến thức nền mà khóa luận dựa vào. Nhóm thứ nhất là bài toán
ghép chuyến theo yêu cầu và kiến trúc giải bài toán đó đang được dùng phổ biến. Nhóm thứ hai
là cách văn liệu hiện hành đo độ bất ổn của lời hứa công bố, cùng với giới hạn của cách đo ấy.
Nhóm thứ ba là khái niệm ổn định lịch trình trong hoạch định sản xuất --- một lĩnh vực đã giải
quyết vấn đề có cấu trúc rất giống nhưng chưa được bắc cầu sang điều phối xe. Nhóm thứ tư là
một phân biệt cổ điển trong luật nghĩa vụ, hóa ra lại mô tả chính xác kết quả cấu trúc mà
chương 3 sẽ chứng minh. Cuối chương là các hệ số định giá độ tin cậy, và một mục ngắn đặt hai cơ
sở tính sổ cái vào từ vựng của điều khiển trực tuyến có nhiễu.

\section{Bài toán ghép chuyến theo yêu cầu}

\subsection{Phát biểu bài toán}

Xét một mạng lưới giao thông biểu diễn bằng đồ thị có hướng $G=(N,A)$, trong đó mỗi cung mang
thời gian di chuyển. Một đội xe $V$ hoạt động trên mạng này; mỗi xe $v \in V$ có vị trí hiện
tại và sức chứa $c_v$. Các yêu cầu đi xe xuất hiện theo thời gian; yêu cầu $r$ được đặc trưng
bởi bộ $\langle o_r, d_r, t_r, q_r \rangle$ gồm điểm đón, điểm trả, thời điểm xuất hiện và số
chỗ cần.

Hệ thống hoạt động theo lô. Cứ sau mỗi chu kỳ $\delta$, toàn bộ các yêu cầu mới tích lũy được
đem gán một lượt cùng với các yêu cầu cũ chưa hoàn thành. Chất lượng dịch vụ được ràng buộc
bởi hai tham số: thời gian chờ tối đa $\Omega$ và độ trễ tối đa $\Delta$ so với chuyến đi
thẳng. Một phương án gán là hợp lệ nếu mọi hành khách được phục vụ đều thỏa hai ràng buộc đó,
và tại mọi thời điểm số khách trên xe không vượt sức chứa.

Bài toán này là dạng động, trực tuyến của Dial-a-Ride Problem, và thuộc lớp NP-khó. Trong thực
tế, các hệ thống không tìm nghiệm tối ưu toàn cục mà giải lại bài toán ở mỗi chu kỳ với trạng
thái hiện hành.

\subsection{Kiến trúc Request--Trip--Vehicle}

Đóng góp kỹ thuật bền vững nhất của Alonso-Mora và cộng sự \cite{alonsomora2017} là tách bài
toán thành hai giai đoạn độc lập. Giai đoạn thứ nhất \emph{sinh} các khả năng: hệ thống dựng
đồ thị ba tầng gồm các yêu cầu, các nhóm yêu cầu có thể phục vụ chung trong một chuyến, và các
cặp chuyến--xe khả thi. Việc tìm các nhóm khả thi được tăng tốc nhờ một tính chất quan trọng:
nếu xe $v$ phục vụ được nhóm $G$ thì nó cũng phục vụ được mọi tập con của $G$, nên có thể mở
rộng nhóm dần theo kích thước và cắt tỉa sớm.

Giai đoạn thứ hai \emph{chọn} phương án bằng một bài toán quy hoạch tuyến tính nguyên. Gọi
$\Gamma$ là tập các chuyến khả thi đã sinh, $x_T \in \{0,1\}$ chỉ việc chuyến $T$ có được thực
hiện hay không, và $z_r \in \{0,1\}$ chỉ việc yêu cầu $r$ bị từ chối. Bài toán có dạng

\begin{equation}
\min_{x,z} \quad \sum_{T \in \Gamma} \mathrm{cost}(T)\, x_T \;+\; \sum_{r \in R} p_{\mathrm{KO}}\, z_r
\end{equation}
\begin{align}
\text{v.đ.k.}\quad & z_r + \sum_{T:\, r \in \mathrm{req}(T)} x_T = 1 && \forall r \in R\\
& \sum_{T:\, \mathrm{veh}(T)=v} x_T \le 1 && \forall v \in V
\end{align}

Ràng buộc (2.2) buộc mỗi yêu cầu hoặc được phục vụ bởi đúng một chuyến, hoặc bị từ chối với
hình phạt $p_{\mathrm{KO}}$; ràng buộc (2.3) buộc mỗi xe nhận không quá một chuyến. Đây là bài
toán đóng gói tập, và nó là khung mà hệ thống của khóa luận cũng sử dụng.

\subsection{Số hạng thiệt hại gây cho khách cũ}

Chi phí của một chuyến trong dòng bài này không chỉ gồm chi phí của những khách mới được phục
vụ. Fielbaum và cộng sự \cite{fielbaum2022scale} viết nó dưới dạng

\begin{equation}
\mathrm{cost}(T,\pi) = \sum_{r \in \mathrm{req}(T)} c_U(r,T,\pi)
\;+\; \sum_{r \in S_{\mathrm{veh}(T)}} \Delta c_U(r,T,\pi)
\;+\; \Delta c_O(\pi)
\end{equation}

trong đó số hạng thứ hai, theo đúng lời các tác giả, là ``chi phí phát sinh thêm cho những hành
khách đang được phục vụ bởi chiếc xe trước khi có phân công này''. Nghĩa là văn liệu \emph{đã}
định giá thiệt hại gây cho khách cũ.

Điểm khác biệt cần nhấn mạnh là số hạng đó có ba đặc điểm. Nó là một đại lượng \emph{biên tại
một thời điểm} --- so kế hoạch mới với kế hoạch cũ ngay lúc đó, không nhớ gì về những lần trước.
Nó được \emph{cộng gộp} trên toàn bộ hành khách bị ảnh hưởng, nên người chịu thiệt nặng nhất bị
pha loãng trong đám đông. Và nó nằm trong \emph{hàm mục tiêu}, tức là một ưu tiên mềm có thể bị
đánh đổi, chứ không phải một ràng buộc. Ba đặc điểm này là lý do vì sao, như chương 4 sẽ cho
thấy, một hành khách vẫn có thể bị dời giờ hẹn tích lũy tới mười bốn phút mà hệ thống không hề
coi đó là bất thường.

\section{Độ bất ổn của lời hứa công bố}

\subsection{Công trình gần nhất}

Fielbaum và Alonso-Mora \cite{fielbaum2020} là công trình gần nhất với khóa luận này, và cần
được trình bày cẩn thận vì phần trùng lặp là đáng kể. Họ đặt đúng vấn đề mà chương 1 đã nêu:
khi hệ thống nhận khách mới, khách cũ bị thiệt bao nhiêu so với điều đã được báo lúc đầu.

Đóng góp của họ gồm một hệ thống định nghĩa và một tập chỉ số. Về định nghĩa, họ phân biệt
nguồn bất ổn \emph{do người dùng gây ra} (khách không có mặt đúng giờ, khách hủy chuyến) với
nguồn \emph{do nhà điều hành gây ra} (hệ thống thay đổi phân công), và phân biệt bất ổn
\emph{một lần} --- xảy ra trong một chuyến cụ thể --- với bất ổn \emph{hàng ngày} --- sự khác
nhau giữa các lần thực hiện cùng một yêu cầu. Về chỉ số, họ định nghĩa

\begin{equation}
\Delta t_w(r) = t_w(r) - t^1_w(r), \qquad
\Delta t_d(r) = t_d(r) - t^1_d(r), \qquad
\Delta D(r) = \Delta t_w(r) + \Delta t_d(r)
\end{equation}

trong đó $t^1_w(r)$ là thời gian chờ \emph{được báo lần đầu tiên}, còn $t_w(r)$ là thời gian
chờ \emph{thực tế}. Các giá trị âm được cắt về không. Họ cũng đếm số lần dự báo tăng lên trong
suốt chuỗi phân công, ký hiệu $NC$.

Kết quả của họ trên Manhattan: hơn một phần ba số yêu cầu không bị từ chối ngay đều gặp ít nhất
một thay đổi, và khi một hành khách bị tăng thời gian chờ thì phần tăng thêm đó lớn tương đương
thời gian chờ trung bình của cả hệ thống. Họ cũng thử ba chế độ kiểm soát và kết luận rằng
kiểm soát chặt hơn thì chất lượng dịch vụ tổng thể xấu đi.

\subsection{Ba giới hạn của cách đo này}

Cách đo trên có ba giới hạn mà khóa luận này nhắm vào.

\textbf{Thứ nhất, nó là hiệu hai đầu chứ không phải độ dài đường đi.} Công thức (2.5) chỉ so
sánh điểm cuối với điểm đầu. Một hành khách được báo 9:00, rồi 9:10, rồi quay về 9:00 sẽ được
chấm $\Delta = 0$, trong khi thực tế người đó đã bị làm phiền hai lần. Chỉ số đếm $NC$ có ghi
nhận chuỗi nhưng lại không mang độ lớn. Nghĩa là văn liệu có \emph{độ lớn nhân với hai đầu} và
có \emph{số lần nhân với chuỗi}, nhưng không có \emph{độ lớn nhân với chuỗi} --- đại lượng mà khóa
luận đặt trần. (Trong hệ của khóa luận, kiểu ``quay về'' như ví dụ trên đến từ trôi dạt hai chiều
chứ không từ quyết định, vì bộ sinh chỉ chèn thêm; mục \ref{muc:dondieu}.)

\textbf{Thứ hai, và quan trọng hơn, thước đo này về mặt toán học không thể trở thành ràng
buộc.} Đại lượng $t_w(r)$ là thời gian chờ thực tế, chỉ biết được sau khi hành khách đã được
đón. Tại thời điểm hệ thống phải quyết định chọn phương án nào, giá trị đó chưa tồn tại. Không
thể chặn ``thực tế trừ lần báo đầu không quá 30 giây'' khi đang chọn phương án, vì chưa biết
thực tế. Đây không phải lựa chọn thiết kế mà là tính chất của công thức: nó chỉ tồn tại ở thì
quá khứ. Cơ chế của khóa luận, ngược lại, định nghĩa trên các đại lượng đều biết được tại thời
điểm quyết định, nên chặn được.

\textbf{Thứ ba, trôi dạt vật lý bị loại khỏi phạm vi thay vì được tách ra.} Các tác giả nêu rõ
trong mục 2 rằng họ không xét những nguồn bất ổn vốn đã có trong giao thông công cộng truyền
thống, chẳng hạn tắc đường hay tai nạn; và trong phần thực nghiệm họ chọn khung giờ 13--15 giờ,
ngoài cao điểm, để tránh đúng hiện tượng đó. Cách xử lý ấy hợp lý trong phạm vi bài của họ,
nhưng nó không dùng được cho giờ cao điểm, và nó khiến câu hỏi ``phần nào do nhà xe, phần nào
do đường'' không thể đặt ra được.

\section{Ổn định lịch trình trong hoạch định sản xuất}

Hiện tượng mà khóa luận nghiên cứu không mới trong khoa học quản trị vận hành; nó chỉ chưa
được chuyển sang lĩnh vực điều phối xe. Trong hoạch định nhu cầu vật tư, từ thập niên 1980 đã
có khái niệm \emph{nervousness}: hệ thống lập lịch chạy theo chân trời cuốn chiếu, mỗi lần chạy
lại sinh ra một lịch mới khác lịch cũ, và sự dao động đó lan xuống toàn bộ chuỗi cung ứng gây
tốn kém.

Điểm đáng chú ý là dòng văn liệu này \emph{đã phân biệt sẵn} hai nguồn dao động: dao động do
thông tin nhu cầu thay đổi --- ngoại sinh --- và dao động do chính quy tắc lập kế hoạch ---
nội sinh. Đó chính là phân biệt mà chương 3 sẽ xây dựng lại cho điều phối xe, với mốc tham
chiếu là vị trí thật của phương tiện thay vì thông tin nhu cầu.

Sridharan, Berry và Udayabhanu \cite{sridharan1987} nghiên cứu việc \emph{đóng băng} một phần
lịch sản xuất chính như biện pháp giảm dao động, và đo đường cong đánh đổi giữa tỉ lệ đóng băng
và sai số chi phí. Kết quả của họ cho thấy chi phí tăng vọt khi tỉ lệ đóng băng vượt khoảng
$80\%$ chân trời. Cơ chế đóng băng ấy tương ứng chính xác với khóa lời hứa trong khóa luận này,
và đường cong đánh đổi của họ tương ứng với đường cong giá của cam kết ở chương 4. Một trong
những quan sát liên ngành mà khóa luận rút ra là hai đường cong này có \emph{hình dạng khác
nhau}, và lý do của sự khác nhau ấy nói lên điều gì đó về cấu trúc của hai bài toán.

Trong một bài đồng hành, nhóm tác giả trên trình bày cả một \emph{họ} thước đo bất ổn, cho phép
đặt trọng số khác nhau lên những thay đổi xảy ra ở các vị trí khác nhau trên chân trời kế hoạch,
và cho phép dùng cả biểu thức tuyến tính lẫn bậc hai. Đây là nền cho phân tích độ vững ở mục
4.5: nếu thứ hạng các cấu hình đổi khi đổi thước đo thì chính điều đó là một phát hiện.

\section{Nghĩa vụ phương tiện và nghĩa vụ kết quả}

Phân biệt giữa hai cơ sở tính sổ cái --- nền của khung bốn họ cổng, kết quả lý thuyết trung tâm
ở chương 3 --- hóa ra trùng với một phân biệt cổ điển trong luật nghĩa vụ, và việc nêu nó ra giúp
kết quả kỹ thuật có chỗ đứng khái niệm rõ ràng hơn.

Luật dân sự phân biệt \emph{nghĩa vụ phương tiện} (obligation de moyens) với \emph{nghĩa vụ kết
quả} (obligation de résultat). Trong nghĩa vụ phương tiện, người có nghĩa vụ không vi phạm nếu
đã hành xử cẩn trọng, kể cả khi kết quả không đạt được; trong nghĩa vụ kết quả, người có nghĩa
vụ chịu trách nhiệm vì không đạt kết quả, bất kể đã cố gắng thế nào. Ví dụ kinh điển: luật sư
chịu nghĩa vụ phương tiện --- cố hết sức nhưng không bảo đảm thắng kiện; thợ lắp thang máy chịu
nghĩa vụ kết quả --- thang chạy hoặc không.

Phân biệt này tồn tại vì một lý do thực dụng mà luật đã kết luận từ lâu: \emph{không thể áp đặt
nghĩa vụ kết quả khi một phần kết quả nằm ngoài tầm kiểm soát của người có nghĩa vụ}, vì khi đó
hợp đồng trở nên bất khả thi. Trong hợp đồng dịch vụ công nghệ thông tin, đúng vấn đề này được
xử lý bằng điều khoản loại trừ: bên cung cấp không chịu trách nhiệm về những sự cố mức dịch vụ
phát sinh từ các sự kiện nằm ngoài tầm kiểm soát hợp lý của mình.

Chương 3 sẽ chứng minh rằng hai cơ sở tính sổ cái cam kết --- tính theo hành vi của nhà điều
hành, hay tính theo kết quả mà khách nhìn thấy --- tương ứng chính xác với hai loại nghĩa vụ
này, và rằng cơ sở kết quả có thể khiến hệ thống bất khả thi đúng theo nghĩa mà luật đã cảnh
báo. Điểm mới nằm ở chỗ: luật và các thỏa thuận mức dịch vụ xử lý ranh giới ``ngoài tầm kiểm
soát hợp lý'' bằng \emph{văn xuôi} --- điều khoản bất khả kháng, tranh cãi sau sự việc --- còn
mốc tham chiếu ngoại sinh trong khóa luận này \emph{tính được} ranh giới đó, theo từng quyết
định và từng hành khách.

\section{Giá trị của độ tin cậy}

Phần cuối của chương ghi lại các hệ số định giá được dùng ở mục \ref{muc:phucloi} để quy đổi kết quả kỹ thuật
sang ngôn ngữ phúc lợi.

Alonso-González và cộng sự \cite{alonsogonzalez2020} ước lượng giá trị thời gian và giá trị độ
tin cậy cho dịch vụ đi chung theo yêu cầu bằng thực nghiệm lựa chọn công bố. Giá trị thời gian
trên xe của họ nằm trong khoảng 7{,}88--10{,}80 EUR mỗi giờ. Đáng chú ý, giá trị độ tin cậy mà
họ đo được --- định nghĩa qua độ lệch chuẩn của thời gian chờ và thời gian đi --- \emph{thấp
hơn} giá trị thời gian tương ứng, chứ không cao hơn. Đây là một điểm mà khóa luận đã phải đính
chính trong quá trình thực hiện: giả định ban đầu rằng độ tin cậy được định giá cao gấp hai tới
ba lần thời gian là sai đối với bối cảnh này, và việc sửa lại đã làm thay đổi đáng kể lập luận
phúc lợi.

Geržinič và cộng sự \cite{gerzinic2023} định giá riêng phần thời gian chờ \emph{bất ngờ}, tức phần
phát sinh ngoài dự kiến ban đầu. Nếu dùng được, đây là hệ số gần nhất với đại lượng mà khóa luận
đo, vì độ dời của một lời hứa đã công bố là thời gian bất ngờ theo nghĩa đó. Mục \ref{muc:phucloi}
dùng khoảng 20--25 EUR mỗi giờ, quy cho nguồn này, để tính điểm hòa vốn; con số ấy khóa luận
\emph{chưa đối chiếu được với bản gốc}, nên mục đó báo cáo song song ngưỡng tính bằng giá trị thời
gian trên xe ở trên, và coi kết quả là một ngưỡng chứ không phải một tuyên bố chiến thắng.

\section{Điều khiển trực tuyến có nhiễu: từ vựng cho hai cơ sở tính}
\label{muc:dieukhien}

Mục này không đưa ra một kết quả mới; nó chỉ đặt hai cơ sở tính sổ cái vào đúng họ bài toán, để
người đọc có sẵn từ vựng. Đây là một \emph{tương ứng cấu trúc}, không phải một kết quả về tỉ số
cạnh tranh đã được chứng minh cho hệ của khóa luận.

Viết lại động học của một lời hứa bằng các ký hiệu sẽ định nghĩa ở chương 3: tại mỗi lần công bố,
lời hứa mới bằng lời hứa cũ cộng phần thay đổi khi giữ nguyên kế hoạch cộng phần do quyết định,
$p_t = v_t + x_t + c_t$, và lời hứa mới trở thành lời hứa cũ của lần sau, $v_{t+1} = p_t$. Đây đúng
là dạng $s_{t+1} = s_t + w_t + u_t$ của \emph{điều khiển trực tuyến với nhiễu} \cite{agarwal2019},
với nhiễu $w_t = x_t$ và tác động điều khiển $u_t = c_t$. Nhiễu ở đây chỉ \emph{gần} ngoại sinh:
$x_t$ gồm cả tiến độ thực của xe, mà tiến độ ấy phụ thuộc tuyến đã chọn, nên về nguyên tắc nó có
thể phụ thuộc yếu vào chính sách (mục \ref{muc:phanra}).

\begin{table}[htbp]
\centering
\caption{Hai cơ sở tính sổ cái trong từ vựng của điều khiển có nhiễu}
\label{bang:dieukhien}
\small
\begin{tabular}{@{}p{3.6cm}p{4.6cm}p{5.4cm}@{}}
\toprule
\textbf{Khóa luận} & \textbf{Điều khiển có nhiễu} & \textbf{Hệ quả}\\
\midrule
Cơ sở hành vi $\sum_t |c_t|$ & phạt \emph{nỗ lực điều khiển} $\sum_t \|u_t\|$ & luôn thỏa được bằng
$u = 0$, tức không làm gì\\
Cơ sở kết quả $\sum_t |x_t + c_t|$ & phạt \emph{dịch chuyển trạng thái}
$\sum_t \|s_{t+1} - s_t\|$, gồm cả nhiễu & có thể vô nghiệm khi riêng nhiễu đã vượt ngân sách\\
\bottomrule
\end{tabular}
\end{table}

Trong tối ưu lồi trực tuyến có làm trơn \cite{chen2018} và trong hệ nhiệm vụ metric
\cite{borodin1992}, chi phí đổi hành động $\sum_t \|a_t - a_{t-1}\|$ chính là tổng biến thiên ---
nhưng ở đó người ra quyết định \emph{đặt trực tiếp} $a_t$, không có trôi dạt ngoại sinh, nên hai cơ
sở tính trùng nhau. Phân biệt hành vi và kết quả chỉ xuất hiện khi trạng thái trôi ngoại sinh, và
lời hứa về giờ giấc đúng là một trạng thái như vậy. Nói gọn: cơ sở hành vi tương ứng với việc chỉ
phạt nỗ lực điều khiển thay vì phạt dịch chuyển trạng thái --- trong điều khiển có nhiễu, chi phí
thường gồm cả hai phần, và cơ sở hành vi chỉ ứng với phần điều khiển; cơ sở kết quả quy nhiễu cho
bộ điều khiển, và đó là lý do nó có thể vô nghiệm.
